\chapter{Repairing a noise obstruction by changing the measurement lag}
\label{ch:pa-retiming}

The noisy example in Chapter~\ref{ch:pa-joint-analog} cannot produce a
positive population gap anywhere in its chosen two-lag quadratic family.
That obstruction concerns a specific set of observation times. Retaining the
same detector and changing those times exposes a negative correlation, which
does admit a dissipation certificate. The new experiment below recovers
386 positive certificates in 400 fresh noisy datasets, compared with zero
for either former fixed contrast, at the same declared evolution time and
with substantially fewer detector readings.

This is a complete example of the research loop: certify why the present
protocol fails, change an identifiable feature of physical access, and test
the resulting prediction on new data. Negative-correlation constraints and
the use of temporal correlations to study irreversibility have established
antecedents~\cite{vlAldousFill2002,ciVSR2024,ciVSRProofs2024}. The contribution
here is the explicit repair, its direct finite-data dissipation bound and its
frozen comparison. It is not a claim that useful timing is new or that every
unknown system has a similarly favorable lag.

\section{One observed lag is sufficient}
Assume the stationary homogeneous Markov and ordinary even-variable reversal
conditions of Theorem~\ref{thm:pa-finite-lag}. Let $Z\in[-1,1]$ be the same
noninvasive memoryless analog detector, with conditional mean $f$ and
independent endpoint noise. Write
\[
 C(\tau)=\mathbb E[Z_0Z_\tau]=\langle f,e^{\tau Q}f\rangle_\pi.
\]
Under a reversible explanation,
\[
 C(\tau)=\|e^{\tau Q/2}f\|_\pi^2\geq0.
\]
This condition includes a nonzero detector mean; no mean subtraction or
zero-lag second-moment estimate is needed for this particular test.

\begin{corollary}[Negative-correlation dissipation certificate]
Set $t=\tau/2$ and $s=0$ in the direct finite-lag theorem. If an already
covered lower bound $g\geq0$ obeys $2g\leq[-C(\tau)]_+$, then
\begin{equation}
 \sigma\geq\frac{2\Phi(2g)}{\tau}\geq\frac{8g}{\tau},\qquad
 \Phi(u)=2\sqrt u\,\operatorname{atanh}\sqrt u,
 \quad 0\leq u<1.
 \label{eq:pa-retiming-bridge}
\end{equation}
The conclusion concerns the same physical model and confidence event.
\end{corollary}
\begin{proof}
The semigroup defect becomes
$-s^2C(0)+2sC(t)-C(2t)=-C(\tau)$.
The two unused coordinates have zero coefficients, so they need not be
measured. The detector mean has range at most two. Substituting $t=\tau/2$
and normalized covered defect $2g$ into the direct theorem gives the result.
\end{proof}

The hidden semigroup identity uses $\tau/2$, but the instrument does not
perform a half-lag measurement. It records only the two endpoints separated
by $\tau$. There are no dedicated zero-lag calibration acquisitions. Bounds
on the detector, its channel stability and memory, the clock and stationary
preparation remain physical calibration requirements.

For $n$ independent stationary pairs, let
$W_i=(1+Z_{i,0}Z_{i,\tau})/2\in[0,1]$. At failure probability $\alpha$,
\begin{equation}
 g_n=\left[\frac12-\overline W
   -\sqrt{\frac{\log(1/\alpha)}{2n}}-e-\delta\right]_+
 \label{eq:pa-retiming-confidence}
\end{equation}
is valid. The quantity $e$ bounds numerical mean error, and $\delta$ bounds
a declared conditional mean discrepancy given the complete preceding record.
The independent stationary-pair design has $\delta=0$. This is Hoeffding's
one-sided bound; it does not assume that an analog score is a Bernoulli
count. The implementation accepts exact product sums, propagates the declared
errors, and applies the same rational enclosure of $\Phi$ as the preceding
chapter. A lag chosen on these same inference records requires a separate
selection correction.

\section{An exact reason to move the lag}
The earlier ring has 20 states, clockwise rate 20 and counterclockwise rate 1.
For an ideal cosine detector, its fundamental-mode correlation has the form
\[
 C(\tau)=A e^{-\gamma\tau}\cos(\omega\tau),\qquad
 \gamma=21\bigl(1-\cos(2\pi/20)\bigr),\quad
 \omega=19\sin(2\pi/20).
\]
The implemented detector is rounded to denominator 4096; the noisy channel
adds an independent fair noise of amplitude $1/2$ to half that cosine.
The exact audit therefore uses the actual rational detector and generator,
not the ideal single-mode approximation.

At the former lags $0.1$ and $0.2$, completing the square and exact
uniformization showed that the greatest analog quadratic gap over every real
$s$ is below $-0.007$. At the newly fixed lag $\tau=0.5$, exact uniformization
instead encloses
\[
 C(0.5)\simeq-0.0732037920< -0.07.
\]
Its population linear dissipation bound is approximately $0.585630$; the
nonlinear bound is $0.600583$. Both remain far below the actual microscopic
rate $19\log20\simeq56.9189$.

The rational calculation expands $e^{\tau Q}$ by uniformization with rate 21.
Each discrete-step correlation has absolute value at most $A$, by stationary
Markov contraction in $L^2(\pi)$. A geometric bound on the positive exponential
tail therefore encloses the omitted correlation contribution, including either
sign. Eighty terms give the reported rigorous interval. An independent
100-digit Fourier calculation checks that interval. This calculation and the
choice of the new lag use the known model after the earlier failed protocol;
the subsequent evaluation uses a new random seed and new held-out records.

\section{Fresh evaluation with equal evolution time}
Before generating the new evaluation bank, the design fixes
$\tau=1/2$, $n=1800$ and $\alpha=1/20$. The new method uses 1,800 stationary
pair preparations, 3,600 detector readings and 900 units of evolution time.
The comparison retains both former fixed values $s=1/2$ and $s=3/4$, each
using the actual joint-analog API and its stronger direct dissipation bridge.
Those methods receive fresh banks of 4,000 stationary single readings and
3,000 pairs at each of lags $0.1$ and $0.2$: 10,000 acquisitions, 16,000
readings and the same 900 units of evolution time.

Thus the retimed design uses 22.5\% as many readings, a factor $40/9$
reduction. Its preparation count also changes. Preparation, reset, detector
operation and processing overhead are unmeasured and are not assigned zero.
The comparison matches declared evolution time; it does not claim equal
total wall time, equal processing energy or a calibrated joule saving.

Four cases each receive 400 new datasets. The driven rates are 20 and 1;
the equilibrium rates are $21/2$ in each direction. Both clean and noisy
versions use the same channels defined above.
\begin{center}
\begin{tabular}{lrrr}
\toprule
Case & Old $s=1/2$ & Old $s=3/4$ & Retimed\\
\midrule
Driven &400/400&400/400&400/400\\
Equilibrium &0/400&0/400&0/400\\
Noisy driven &0/400&0/400&386/400\\
Noisy equilibrium &0/400&0/400&0/400\\
\bottomrule
\end{tabular}
\end{center}
The noisy detection fraction is 96.5\%, with a nominal Wilson 95\% Monte
Carlo interval of 94.21--97.90\%. Its mean nonlinear lower bound is
$0.124947$, only 0.2195\% of the actual rate. Retiming restores a usable
certificate in this example; it does not reconstruct the microscopic
dissipation accurately. For the clean driven case, the mean bound is
$2.05450$, compared with $0.621640$ and $0.900411$ for the two earlier fixed
designs. None of the 1,600 datasets produces a bound above its model's true
rate. That simulation observation is additional checking, not a replacement
for the confidence proof.

All three methods have individual 95\% certificates. Their unadjusted maximum
is not advertised as another 95\% certificate. The archive preserves every
dataset's exact sufficient statistics and the first raw bank of each case.
Independent code reconstructs every reported rate and checks the raw-bank
summaries, physical counts and population intervals.

\section{What this repair does and does not identify}
The timing change avoids the conservative $B-A$ calibration penalty because
the zero-lag coefficient is zero. Independent detector noise still reduces
the conditional mean signal and hence the accessible correlation. It does
not vanish from the experiment, and many driven systems or available lag
ranges may have no statistically useful negative correlation.

The ordinary-reversal and detector assumptions also remain necessary. An
equilibrium one-state target coupled to an unobserved driven sensor ring can
produce exactly the same recorded law as the noisy driven example. The full
target-plus-sensor system dissipates, while the target alone need not.
Assigning the inferred lower bound to the target requires excluding such
sensor memory or including it in the physical system being bounded.
Equilibrium systems with odd hidden variables, such as inertial dynamics,
likewise require the correct reversal operation and a separate argument.

The result establishes a concrete change in experimental access: one protocol
family provably has no positive population gap, while a redesigned timing
choice provides a finite-data dissipation certificate with fewer readings at
the same evolution budget. It is a model-assisted, reproducible repair under
stated assumptions, rather than a universal rule for selecting lags without
prior information about a system.
